\iffalse
Literature review entry for ELF: Embedded Language Flows (Hu et al. 2026, arxiv:2605.10938).
Written per the write-latex skill rules:
  - no fontspec
  - \section* (unnumbered) by default
  - pdflatex-clean
  - comments only inside \iffalse ... \fi
\fi

\documentclass[11pt]{article}

\usepackage[margin=1in]{geometry}
\usepackage{amsmath, amssymb, amsthm}
\usepackage{mathtools}
\usepackage{natbib}
\usepackage{hyperref}
\usepackage{xcolor}

\hypersetup{
  colorlinks=true,
  linkcolor=black,
  citecolor=blue!50!black,
  urlcolor=blue!50!black
}

\title{Embedded Language Flows: A Minimalist Bridge Between Image-Domain Flow Matching and Diffusion Language Modeling --- A Review of \citet{hu2026elf}}
\author{Literature review entry}
\date{Studied: 2026-05-14 \\ \texttt{arXiv:2605.10938}}

\begin{document}
\maketitle

\begin{abstract}
\citet{hu2026elf} propose Embedded Language Flows (ELF), a class of continuous diffusion language models (DLMs) based on continuous-time Flow Matching \citep{lipman2022fm, liu2022rectified} operating in a frozen pretrained embedding space. Unlike prior continuous DLMs that supervise discretization at every timestep, ELF keeps the entire sampling trajectory in continuous embedding space and only discretizes at $t = 1$ via the weight-tied unembedding matrix, eliminating the need for a separately trained decoder. The formulation is structurally identical to image-domain latent diffusion \citep{rombach2022ldm}, which lets the authors directly transfer classifier-free guidance \citep{ho2022cfg} and other established image-domain techniques. With $\sim 105$M parameters, $\sim 10\times$ fewer training tokens than baselines, and fewer sampling steps, ELF outperforms leading discrete DLMs such as MDLM \citep{sahoo2024mdlm} and SEDD \citep{lou2024sedd}, as well as continuous DLMs descended from Diffusion-LM \citep{li2022diffusionlm}. The result challenges the dominant framing that language requires fundamentally discrete diffusion machinery and offers a cleaner template for cross-modal generative modeling.
\end{abstract}

\section*{Background and Motivation}

Diffusion and flow-based generative models \citep{ho2020ddpm, lipman2022fm, liu2022rectified, albergo2025interpolants} have become the dominant paradigm for continuous data (images, video, audio). Their extension to language has split into two camps. Discrete DLMs \citep{austin2021d3pm, sahoo2024mdlm, lou2024sedd} apply diffusion directly over the categorical token vocabulary, using absorbing-state or score-ratio formulations. Continuous DLMs \citep{li2022diffusionlm} lift discrete tokens into a continuous embedding space, run diffusion there, and discretize back. Continuous DLMs have historically underperformed discrete DLMs because they require (i) per-step cross-entropy supervision to keep the trajectory near valid tokens and (ii) a separately trained decoder that maps embedding-space samples back to text. Both constraints degrade flow flexibility and add architectural overhead.

The paper under review removes both constraints.

\section*{Technical Approach}
\label{sec:method}

Let $E: \mathcal{V} \to \mathbb{R}^d$ denote a token embedding map, $z_1 = E(x)$ the clean embedding tensor for sequence $x \in \mathcal{V}^L$, and $z_0 \sim \mathcal{N}(0, I)$. ELF uses a linear interpolation path
\begin{equation}
z_t \;=\; (1 - t) z_0 + t \, z_1, \qquad t \in [0, 1].
\label{eq:path}
\end{equation}

\paragraph{Flow matching objective.}
A neural network $v_\theta(z_t, t)$ is trained to predict the target velocity $z_1 - z_0$ via
\begin{equation}
\mathcal{L}_{\text{FM}}(\theta) \;=\; \mathbb{E}_{x, z_0, t}\Big[\, \big\| v_\theta(z_t, t) - (z_1 - z_0) \big\|^2 \,\Big].
\label{eq:fm}
\end{equation}
Equivalently the network may predict the clean embedding $\hat{x}_\theta(z_t, t) = \mathbb{E}[z_1 \mid z_t]$ (the $x$-prediction parameterisation used by the authors).

\paragraph{Final-step discretization via weight-tied unembedding.}
At inference, $z_0$ is integrated forward to $t = 1$. ELF's key observation is that the same encoder that produced $z_1 = E(x) = W_E x$ can be re-used at the final step as the decoder: discrete tokens are recovered via
\begin{equation}
\pi_\theta(z_t \to 1) \;=\; \mathrm{softmax}\big( W_E \cdot \hat{x}_\theta(z_{t \to 1}, 1) \big),
\label{eq:decode}
\end{equation}
where $W_E \in \mathbb{R}^{|\mathcal{V}| \times d}$ is shared between the encoder and the unembedding head. No separate decoder is trained. Cross-entropy supervision is applied \emph{only} at $t = 1$, with weight $\alpha$, alongside the flow matching loss:
\begin{equation}
\mathcal{L}(\theta) \;=\; \mathcal{L}_{\text{FM}}(\theta) + \alpha \, \mathcal{L}_{\text{CE}}(\theta).
\end{equation}

\paragraph{Classifier-free guidance with self-conditioning.}
ELF inherits CFG \citep{ho2022cfg} directly. The conditioning signal $c$ is the model's own intermediate prediction (self-conditioning); CFG combines conditional and unconditional velocity fields as
\begin{equation}
v_{\text{cfg}}(z_t \mid c) \;=\; \omega \, v_\theta(z_t \mid c) + (1 - \omega) \, v_\theta(z_t \mid \emptyset).
\label{eq:cfg}
\end{equation}
To avoid the $2\times$ inference cost of naive CFG, the authors adopt \emph{training-time CFG} (cf.\ image-domain analogues): train the network to output $v_\text{cfg}$ directly in one forward pass.

\paragraph{Embedding choice.}
The encoder $E$ is studied in three regimes: pretrained contextual (transformer input embeddings), non-contextual learned (lookup table), and frozen random projection. Pretrained contextual embeddings dominate the ablation, suggesting the well-structured embedding manifold provides much of the lift.

\section*{Key Results}

ELF is evaluated on three setups. (i) Unconditional generation on OpenWebText, with generative perplexity under a pretrained GPT-2 Large evaluator and unigram entropy as a diversity proxy. (ii) Conditional generation on WMT14 De$\to$En translation. (iii) Conditional generation on summarization.

The headline empirical findings: with a $\sim 105$M parameter model trained on $\sim 10\times$ fewer tokens than its $\sim 170$M-parameter competitors, ELF achieves lower generative perplexity at fewer sampling steps than MDLM \citep{sahoo2024mdlm}, SEDD \citep{lou2024sedd}, and prior continuous DLMs. CFG sweeps trace a quality--diversity Pareto front (higher $\omega$ lowers perplexity, lowers entropy). Ablations confirm: pretrained contextual embeddings dominate; per-step CE supervision (the prior-art approach) underperforms ELF's at-final-step-only formulation; training-time CFG matches inference-time CFG quality at half the cost.

\section*{Position in the Literature}

ELF is a structural simplification of prior continuous DLMs. Where \citet{li2022diffusionlm} apply per-step categorical supervision to keep the trajectory close to valid tokens, ELF lets the trajectory run free and discretizes only at the final step. Where latent diffusion \citep{rombach2022ldm} trains a separate decoder, ELF reuses the encoder via weight tying. Where discrete DLMs \citep{austin2021d3pm, sahoo2024mdlm, lou2024sedd} avoid the embedding-space question entirely by formulating diffusion over categoricals, ELF returns to the embedding-space formulation and shows it can be made competitive with minimal machinery.

The broader thesis: language can be modeled with the same flow-matching toolkit developed for images, provided the discrete--continuous interface is handled at exactly one boundary (the final time step) rather than diffusely across the trajectory. This is a unification claim in the sense that the empirical machinery is shared with image-domain flow matching \citep{lipman2022fm, liu2022rectified}.

\section*{Limitations and Open Questions}

The empirical wins come with caveats. The pretrained contextual embedding is doing significant work; the authors do not quantitatively decompose how much of the gap over discrete baselines is due to the flow matching formulation versus the embedding manifold. Generative perplexity under GPT-2 Large is the headline quality metric, which biases the evaluation toward models whose output distribution resembles GPT-2's; unigram entropy is a weak diversity proxy that does not catch $n$-gram-level repetition or mode collapse. No comparison is reported against a similarly-sized autoregressive transformer trained on the same OpenWebText setup with matched token budget --- the salient missing control.

The mathematical setup leaves the choice of linear interpolation path unexamined: rectified-flow-style straight-line paths are convenient but not obviously optimal for embedding spaces with non-trivial geometry. The CFG scale $\omega$ is swept empirically; no analysis is offered for when the quality--diversity Pareto front saturates or inverts.

\section*{Sandbox Probe}
A minimal CPU-runnable PyTorch experiment was built to probe the weight-tied-unembedding mechanism in isolation. A tiny MLP is trained to predict $\hat{x}_\theta$ along the linear flow path on a synthetic vocabulary of $64$ tokens with $32$-dimensional random embeddings; Euler integration with varying step counts is followed by $\arg\max W_E z$ decoding. Source: \url{https://github.com/pleyva2004/elf-embedded-language-flows} (sandbox subdirectory).

\bibliographystyle{plainnat}
\bibliography{references}

\end{document}
